% Part: proof-theory
% Chapter: natural-deduction
% Section: quantifiers

\documentclass[../../../include/open-logic-section]{subfiles}

\begin{document}

\olfileid{pt}{nat}{qua}

\olsection{নিয়মিত \usetoken{P}{proof} ও প্রতিস্থাপন}

\Elim{\lexists} ও~\Intro{\lforall} বিধিতে আইগেনচলের শর্ত
থাকায় পরিমাণসূচকের বিধিগুলি কিছু জটিলতা আনে। এই
অনুমানগুলিতে !!{constant}~$c$ আর কোথাও ব্যবহার না করলে
বেশির ভাগ জটিলতা সামলানো যায়। তাই একটি সংজ্ঞা দিই।

\begin{defn}
স্বাভাবিক নিষ্পাদনের !!{proof}~$\delta$-কে \emph{নিয়মিত}
বলব যদি
\begin{enumerate}
  \item কোনো আইগেনচল~$a$ একাধিক \Elim{\lexists} বা~\Intro{\lforall}
  অনুমানের আইগেনচল হিসেবে ব্যবহৃত না হয়; এবং
  \item $a$ যদি কোনো \Elim{\lexists} বা~\Intro{\lforall}
  অনুমানের আইগেনচল হয়, তাহলে যথাক্রমে গৌণ পূর্বধারণা
  বা একমাত্র পূর্বধারণায় পৌঁছানো অব্যবহিত উপ-!!{proof}-এর
  ভিতরেই শুধু সেটি ঘটে।
\end{enumerate}
\end{defn}

\begin{lem}\ollabel{lem:subst-regular} ধরি, $\delta$ নিয়মিত এবং
তার অন্তিম সূত্র~$!C$; !!a{constant}~$c$ প্রমাণ~$\delta$-তে
আইগেনচল হিসেবে ব্যবহৃত নয়; আর পদ~$t$-তে $\delta$-র
কোনো আইগেনচল নেই এবং পদটি বদ্ধ। তাহলে $\delta$-তে $c$-র প্রতিটি
সংঘটন $t$ দিয়ে বদলে পাওয়া $\Subst{\delta}{t}{c}$-ও
নিয়মিত !!{proof}। $\Subst{\delta}{t}{c}$-র অন্তিম-!!{formula} হলো
$\Subst{!C}{t}{c}$। $!A^x$ যদি $\delta$-র খোলা অনুমান হয়,
তাহলে $\Subst{!A}{t}{c}^x$ হলো
$\Subst{\delta}{t}{c}$-র খোলা অনুমান।
% BN-SRC-874 TARGET-CORRECTION-BEGIN
\par\smallskip\noindent\textnormal{\small\textbf{পাঠ-সংশোধন (BN-SRC-874):} এই অধ্যায়ের বিধিতে প্রতিস্থাপনের পদ বদ্ধ হতে হয়।
উপপাদ্যে সেই শর্ত যোগ করা হয়েছে: চলকযুক্ত পদে ধ্রুবক বদলালে
অস্তিত্ব-ভূমিকার বদ্ধ-পদের শর্ত নষ্ট হতে পারে।}
% BN-SRC-874 TARGET-CORRECTION-END
\end{lem}

\begin{proof}
  $\pheight{\delta}$-র উপর আরোহ করি।
  $\pheight{\delta} = 0$ হলে প্রমাণটি শুধু
  অনুমান~$!A^x$; তখন $\Subst{\delta}{t}{c}$ হলো
  $\Subst{!A}{t}{c}^x$।

  $\pheight{\delta} > 0$ হলে $\delta$-র শেষ অনুমানবিধি
  অনুযায়ী ক্ষেত্রগুলি আলাদা করি। বচনীয় সংযোজকের
  বিধিক্ষেত্রগুলি সহজ; অনুশীলন হিসেবে থাক। এখানে
  $\delta$ পরিমাণসূচকের বিধিতে শেষ হয় এমন ক্ষেত্র দেখি।
  \begin{enumerate}
    \item শেষ বিধি $\Intro{\lforall}$ হলে $\delta$-র রূপ
    \[
    \AxiomC{}
    \RightLabel{$\delta'$}
    \DeduceC{$!A(a,c)$}
    \RightLabel{\Intro{\lforall}}
    \UnaryInfC{$\lforall[x][!A(x,c)]$}
    \DisplayProof
    \]
    আরোহের অনুমানে $\Subst{\delta'}{t}{c}$ হলো
    $!A(a,t)$-র নিয়মিত !!{proof}। $a$ প্রমাণ~$\delta$-র
    আইগেনচল বলে পদ~$t$-তে $a$ ঘটে না। অতএব
    $\Subst{\delta}{t}{c}$-র শেষ অনুমানবিধি,
    \[
    \AxiomC{}
    \RightLabel{$\Subst{\delta'}{t}{c}$}
    \DeduceC{$!A(a,t)$}
    \RightLabel{\Intro{\lforall}}
    \UnaryInfC{$\lforall[x][!A(x,t)]$}
    \DisplayProof
    \]
    সঠিক: $t$-তে $a$ না থাকায় সিদ্ধান্ত বা কোনো
    খোলা অনুমানেও $a$ নেই।
    \item শেষ বিধি $\Elim{\lforall}$ হলে
    $\delta$-র রূপ
    \[
    \AxiomC{}
    \RightLabel{$\delta'$}
    \DeduceC{$\lforall[x][!A(x,c)]$}
    \RightLabel{\Elim{\lforall}}
    \UnaryInfC{$!A(s(c),c)$}
    \DisplayProof
    \]
    আরোহের অনুমানে $\Subst{\delta'}{t}{c}$ হলো
    $\lforall[x][!A(x,t)]$-র নিয়মিত !!{proof}।
    (সিদ্ধান্তের পদ~$s(c)$-তে $c$ থাকতেও পারে,
    না-ও পারে।) $\Subst{\delta}{t}{c}$-র শেষ অনুমানবিধি,
    \[
    \AxiomC{}
    \RightLabel{$\Subst{\delta'}{t}{c}$}
    \DeduceC{$\lforall[x][!A(x,t)]$}
    \RightLabel{\Elim{\lforall}}
    \UnaryInfC{$!A(s(t),t)$}
    \DisplayProof
    \]
    সঠিক, এবং $!A(s(t),t) \ident \Subst{!A(s(c),c)}{t}{c}$।
    \Intro{\lexists}-এর ক্ষেত্রও অনুরূপ।
  \item শেষ বিধি $\Elim{\lexists}$ হলে $\delta$-র রূপ
    \[
    \AxiomC{}
    \RightLabel{$\delta_1'$}
    \DeduceC{$\lexists[x][!A(x,c)]$}
    \AxiomC{$\Discharge{!A(a,c)}{x}$}
    \RightLabel{$\delta_2'$}
    \DeduceC{$!C$}
    \DischargeRule{\Elim{\lexists}}{x}
    \BinaryInfC{$!C$}
    \DisplayProof
    \]
    আরোহের অনুমানে $\Subst{\delta_1'}{t}{c}$ ও
    $\Subst{\delta_2'}{t}{c}$ নিয়মিত !!{proof}s:
    প্রথমটির সিদ্ধান্ত $\lexists[x][!A(x,t)]$;
    দ্বিতীয়টির সিদ্ধান্ত $\Subst{!C}{t}{c}$,
    খোলা অনুমান~$!A(a,t)$ থেকে। $a$ প্রমাণ~$\delta$-র
    আইগেনচল বলে পদ~$t$-তে $a$ নেই। অতএব
    $\Subst{\delta}{t}{c}$-র শেষ অনুমানবিধি,
    \[
    \AxiomC{}
    \RightLabel{$\Subst{\delta_1'}{t}{c}$}
    \DeduceC{$\lexists[x][!A(x,t)]$}
    \AxiomC{$\Discharge{!A(a,t)}{x}$}
    \RightLabel{$\Subst{\delta_2'}{t}{c}$}
    \DeduceC{$\Subst{!C}{t}{c}$}
    \DischargeRule{\Elim{\lexists}}{x}
    \BinaryInfC{$\Subst{!C}{t}{c}$}
    \DisplayProof
    \]
    সঠিক: $\Subst{!A(x,t)}{a}{x} \ident \Subst{!A(a,c)}{t}{c}$;
    সিদ্ধান্ত~$\Subst{!C}{t}{c}$ বা অন্য কোনো খোলা
    অনুমানে $a$ নেই।
    % BN-SRC-544 TARGET-CORRECTION-BEGIN
    \par\smallskip\noindent\textnormal{\small\textbf{উৎস-সংশোধন (BN-SRC-544):} অস্তিত্ব-নিষ্কাশনের আরোহফলে $c$-র স্থলে $t$ বসায় প্রধান পূর্বধারণা ও সিদ্ধান্তে সেই প্রতিস্থাপিত রূপ দেখানো হয়েছে।}
    % BN-SRC-544 TARGET-CORRECTION-END
    % BN-SRC-545 TARGET-CORRECTION-BEGIN
    \par\smallskip\noindent\textnormal{\small\textbf{উৎস-সংশোধন (BN-SRC-545):} উৎসে এই ক্ষেত্রের দ্বিতীয় প্রমাণবৃক্ষ ভুলে সার্বিক-ভূমিকার পুনরাবৃত্তি; এখানে অস্তিত্ব-নিষ্কাশনের দুই পূর্বধারণা ও অব্যাহতি-চিহ্ন দেখানো হয়েছে।}
    % BN-SRC-545 TARGET-CORRECTION-END
  \end{enumerate}
\end{proof}

\begin{prop}\ollabel{prop:regular}
$\delta$ যদি \Log{N1c} বা \Log{N1i}-তে !!a{proof} হয়,
তাহলে একই গঠনবিশিষ্ট (বিশেষত, একই !!{height}-এর)
নিয়মিত !!{proof}~$\delta'$ আছে, যা $\delta$ থেকে
শুধু আইগেনচলগুলি প্রতিস্থাপন করায় আলাদা।
\end{prop}

\begin{proof}
শুধু এই প্রমাণের জন্য, $\delta$-র কোনো আইগেনচল-বিধিকে
\emph{পরিচ্ছন্ন} বলি যদি তার আইগেনচল অন্য কোনো অনুমানবিধির
আইগেনচল না হয় এবং ওই বিধির অব্যবহিত উপ-!!{proof}-র
বাইরে $\delta$-তে কোথাও না ঘটে; অন্যথায় তাকে
\emph{অপরিচ্ছন্ন} বলি। $\delta$-তে অপরিচ্ছন্ন আইগেনচল-বিধির
সংখ্যা~$n$ ধরি। $n$-র উপর আরোহে দেখাই, $\delta$-কে
প্রয়োজনীয় নিয়মিত !!{proof}~$\delta'$-তে বদলানো যায়।
$n=0$ হলে $\delta$-র সব আইগেনচল-বিধিই পরিচ্ছন্ন; তাই $\delta$
নিয়মিত এবং আর কিছু প্রমাণের নেই। অন্যথায় সর্বোচ্চে
থাকা কোনো অপরিচ্ছন্ন আইগেনচল-বিধি বাছি, ধরি
\Intro{\lforall}। তাতে শেষ হওয়া $\delta$-র
উপ-!!{proof}~$\delta_1$-এর রূপ
    \[
    \AxiomC{}
    \RightLabel{$\delta_2$}
    \DeduceC{$!A(c)$}
    \RightLabel{\Intro{\lforall}}
    \UnaryInfC{$\lforall[x][!A(x)]$}
    \DisplayProof
    \]
$c$ আইগেনচল বলে $\lforall[x][!A(x)]$ বা কোনো
খোলা অনুমানে সেটি ঘটে না। উপপ্রমাণ~$\delta_2$
নিয়মিত: তার উপরের সব আইগেনচল-বিধি পরিচ্ছন্ন,
যদিও সেখানে এমন বিধি থাকতেই পারে। প্রমাণ~$\delta$-তে
ঘটে না এমন একটি !!a{constant}~$a$ নিই।
\olref{lem:subst-regular} অনুযায়ী
$\Subst{\delta_2}{a}{c}$ হলো $!A(a)$-র নিয়মিত প্রমাণ।
আইগেনচল-শর্তে $\delta_2$-র কোনো খোলা অনুমানে $c$
থাকে না; অতএব $\Subst{\delta_2}{a}{c}$-র কোনো খোলা
অনুমানেও $a$ নেই। ফলে~$\Intro{\lforall}$ প্রয়োগ করি।
$\delta$-তে $\delta_1$-র বদলে প্রমাণ~$\delta_1'$ বসালে,
    \[
    \AxiomC{}
    \RightLabel{$\Subst{\delta_2}{a}{c}$}
    \DeduceC{$!A(a)$}
    \RightLabel{\Intro{\lforall}}
    \UnaryInfC{$\lforall[x][!A(x)]$}
    \DisplayProof
    \]
একটি অপরিচ্ছন্ন আইগেনচল-বিধি পরিচ্ছন্ন বিধিতে বদলেছে;
শুধু আইগেনচল~$c$-র জায়গায়~$a$ এসেছে। ফলিত
প্রমাণে অপরিচ্ছন্ন বিধি সর্বাধিক $n-1$টি; আরোহের অনুমানে ফল মেলে।
% BN-SRC-546 TARGET-CORRECTION-BEGIN
\par\smallskip\noindent\textnormal{\small\textbf{উৎস-সংশোধন (BN-SRC-546):} আরোহে একটি অপরিচ্ছন্ন বিধিই বেছে নেওয়া দরকার; তার উপরে পরিচ্ছন্ন আইগেনচল-বিধি থাকতে পারে, কিন্তু উপপ্রমাণ তবু নিয়মিত। উৎসের ``সর্বোচ্চ আইগেনচল-বিধি'' ও ``উপরে কোনো আইগেনচল-বিধি নেই'' বক্তব্য সেই নিশ্চয়তা দেয় না।}
% BN-SRC-546 TARGET-CORRECTION-END
% BN-SRC-648 TARGET-CORRECTION-BEGIN
\par\smallskip\noindent\textnormal{\small\textbf{উৎস-সংশোধন (BN-SRC-648):} একটি অপরিচ্ছন্ন আইগেনচলের নাম বদলালে অন্য কোনো অপরিচ্ছন্ন বিধিও পরিচ্ছন্ন হতে পারে; তাই ফলিত প্রমাণে অপরিচ্ছন্ন বিধির সংখ্যা ঠিক \(n-1\) নয়, সর্বাধিক \(n-1\)। আরোহের জন্য এই কঠোর হ্রাসই যথেষ্ট।}
% BN-SRC-648 TARGET-CORRECTION-END
\end{proof}

\begin{prob}
\olref{prop:regular}-এর প্রমাণে সর্বোচ্চ অপরিচ্ছন্ন আইগেনচল-বিধি
যখন~\Elim{\lexists}, সেই ক্ষেত্রটি বর্ণনা করুন।
\end{prob}

এইমাত্র প্রমাণিত প্রস্তাবনা আমরা আলাদা করে আর প্রয়োগ
করব না; বিবেচিত সব !!{proof}s-কে নিয়মিত ধরব—তা না
হলে আইগেনচল বদলে নিয়মিত করা যায়।

\Olref{lem:subst-regular} ও \olref{prop:regular}
\Log{N2c} এবং~\Log{N2i}-তেও প্রযোজ্য। প্রমাণগুলি
অনুশীলন হিসেবে থাক।

\end{document}
